\begin{abstract}
We present a lightweight video face recognition system designed for edge deployment. The core encoder is a MobileNetV4-Conv-Medium student (9.58M parameters, 228M MACs) distilled from a frozen MagFace iResNet-100 teacher (65.2M parameters, 12.15G MACs), achieving a 6.8$\times$ parameter reduction and 53.3$\times$ MACs reduction. The student is trained with a combination of MagFace classification loss and embedding-level knowledge distillation. We study four independently trained variants: a clean KD baseline (V1), a spatial KD variant with aggressive occlusion curriculum (V2), an SWA-stabilized robust variant (V3), and a spatial-KD variant with gated masking curriculum (V4). The clean KD student achieves 87.98\% TAR at FAR=$10^{-4}$ on IJB-B and 90.65\% on IJB-C, within 1.44\% of MobileFaceNet (W600K) on IJB-B and within 1.07\% on IJB-C, despite using 7$\times$ fewer training identities. V4/best achieves the highest LFW accuracy among all students (99.58\%), confirming that spatial KD during clean epochs improves frontal-face discriminability; its partial masking curriculum (epoch 23 of 40) also provides intermediate occlusion robustness between V1 and V3. V3/SWA substantially reduces degradation under synthetic lower-face masking (TAR drop reduced by up to 69\% relative to MBF on CPLFW) while maintaining strong verification performance. On the RMFRD real masked-face dataset (RetinaFace-aligned), V3/SWA (84.18\% AUC, 23.03\% Rank-1) outperforms both V1/best (82.66\%, 22.01\%) and V4/best (82.07\%, 20.67\%), confirming that the full occlusion curriculum and SWA are needed for maximal real-mask robustness. On MFR2 (53 identities), all student variants outperform MBF on verification AUC (V3/SWA: $+4.47\%$), with V4/best achieving the highest student Rank-1 identification (69.01\%) attributed to its superior frontal-face embedding quality. The distilled encoder is integrated into a complete runtime pipeline with YOLO-based face detection, 5-point alignment, multi-object tracking, liveness filtering, MagFace magnitude-based quality gating, magnitude-weighted track-level pooling, and FAISS-based identity retrieval. A template pooling ablation on IJB-B/C confirms that magnitude-weighted pooling consistently outperforms simple mean pooling at strict FAR operating points. Limitations include the absence of standardized open-set 1:N video evaluation and a fully controlled PAD ablation, which are left for future work.
\end{abstract}
